\begin{theorem}[Integer balanced contextual opposite-sign zero collapse]
\label{thm:int-balanced-contextual-opposite-sign-zero-collapse}
For all $a,b,m,n:\mathsf{BHist}$, if $\mathsf{Unary}(a)$,
$\mathsf{Unary}(b)$, $\mathsf{Unary}(m)$, and $\mathsf{Unary}(n)$, then
the balanced contextual opposite-sign pair collapses exactly at the two zero
counters:
$$
\mathsf{IntPairClassifier}
((\mathsf{append}(a,m),\mathsf{append}(\emp,b)),
(\mathsf{append}(a,\emp),\mathsf{append}(n,b)))
\Longleftrightarrow
\hsame(m,\emp)\land\hsame(n,\emp).
$$
The ten-field fit is the following. The traditional target is contextual
opposite-sign integer zero collapse. The BEDC source consists of
$\mathsf{IntPairClassifier}$, $\mathsf{Unary}$, $\mathsf{append}$, $\emp$,
balanced context exactness, and canonical opposite-sign zero collapse. The
carrier is carried by unary difference pairs. The classifier is
cross-addition sameness. Stability is removal of the shared positive and
negative contexts. The ledger records hidden common contexts $a,b$ together
with zero counters. The core theorem is this balanced contextual zero-collapse
block. Exactness is the displayed iff. The standard bridge is ordinary
cancellation of equal added contexts before the fact that a unary sum is zero
only when both summands are zero. The status is $\mathsf{ProofTarget}$.
\end{theorem}
\begin{proof}
For the forward direction, assume the contextual classifier. Apply
Theorem~\ref{thm:int-pair-classifier-balanced-append-context-exactness} with
$p=m$, $n=\emp$, $q=\emp$, and $m=n$. The unary hypotheses for $a$ and $b$
remove the balanced positive and negative contexts and reduce the classifier to
$$
\mathsf{IntPairClassifier}((m,\emp),(\emp,n)).
$$

Theorem~\ref{thm:int-canonical-opposite-sign-zero-collapse}, applied with the
unary hypotheses for $m$ and $n$, converts this canonical opposite-sign
classifier into $\hsame(m,\emp)\land\hsame(n,\emp)$.

For the reverse direction, assume
$\hsame(m,\emp)\land\hsame(n,\emp)$. The canonical
opposite-sign zero-collapse theorem first gives
$\mathsf{IntPairClassifier}((m,\emp),(\emp,n))$ from the unary hypotheses for
$m$ and $n$.

Apply Theorem~\ref{thm:int-pair-classifier-balanced-append-context-exactness}
again with $p=m$, $n=\emp$, $q=\emp$, and $m=n$. Reading this equivalence in
the reverse direction, the unary hypotheses for $a$ and $b$ transport the
canonical classifier into the displayed balanced contextual classifier.
\end{proof}

\begin{theorem}[Integer contextual opposite-sign zero collapse]
\label{thm:int-contextual-opposite-sign-zero-collapse}
For all $a,m,n:\mathsf{BHist}$, if $\mathsf{Unary}(a)$,
$\mathsf{Unary}(m)$, and $\mathsf{Unary}(n)$, then the visible-context
opposite-sign canonical pair collapses exactly at the shared zero context and
zero payloads:
$$
\mathsf{IntPairClassifier}((\mathsf{append}(a,m),\emp),(\emp,\mathsf{append}(n,a)))
\Longleftrightarrow
\hsame(a,\emp)\land\hsame(m,\emp)\land\hsame(n,\emp).
$$
The ten-field fit is the following. The traditional target is contextual
opposite-sign signed-unary zero detection. The BEDC source consists of
$\mathsf{IntPairClassifier}$, $\mathsf{Unary}$, $\mathsf{append}$, $\emp$, and
$\hsame$. The carrier is unary difference pairs. The classifier is
cross-addition sameness. Stability is append-empty inversion. The ledger
records the shared context together with the two payloads. The core theorem is
this contextual zero-collapse block. Exactness is the displayed iff. The
standard bridge is the ordinary fact that $a+m+n+a=0$ for unary summands forces
all displayed summands to be zero. The status is $\mathsf{ProofTarget}$.
\end{theorem}
\begin{proof}
Assume first the three unary hypotheses and the displayed classifier. Unary
append closure makes $\mathsf{append}(a,m)$ and $\mathsf{append}(n,a)$ unary, so
Theorem~\ref{thm:int-canonical-opposite-sign-zero-collapse} applies to the
canonical pair built from these two appended counters.

The forward half of that theorem gives
$\hsame(\mathsf{append}(a,m),\emp)$ and
$\hsame(\mathsf{append}(n,a),\emp)$. Empty readback converts these into two
empty append equations, and append-empty inversion yields
$a=\emp$, $m=\emp$, and $n=\emp$ after reading the shared context from both
equations. Transport along these empty equations gives
$\hsame(a,\emp)\land\hsame(m,\emp)\land\hsame(n,\emp)$.

Conversely, assume
$\hsame(a,\emp)\land\hsame(m,\emp)\land\hsame(n,\emp)$. Empty readback turns the
three readbacks into empty equations for $a$, $m$, and $n$. The append unit laws
therefore identify both contextual counters
$\mathsf{append}(a,m)$ and $\mathsf{append}(n,a)$ with $\emp$.

The reverse half of Theorem~\ref{thm:int-canonical-opposite-sign-zero-collapse}
now applies to these two unary appended counters. Its required empty readbacks
are obtained by the preceding transport to $\emp$, and the resulting canonical
classifier is exactly
$\mathsf{IntPairClassifier}((\mathsf{append}(a,m),\emp),(\emp,\mathsf{append}(n,a)))$.
\end{proof}

\begin{theorem}[Integer history classifier positive append context iff]
\label{thm:int-history-classifier-positive-append-context-iff}
For all $a,m,n:\mathsf{BHist}$, if $\mathsf{Unary}(a)$,
$\mathsf{Unary}(m)$, and $\mathsf{Unary}(n)$, then the public same-sign
integer history classifier is invariant under adding the same visible unary
prefix to both magnitudes:
$$
\mathsf{IntHistoryClassifier}(\Ezero(\mathsf{append}(a,m)),\Ezero(\mathsf{append}(a,n)))
\Longleftrightarrow
\mathsf{IntHistoryClassifier}(\Ezero(m),\Ezero(n)).
$$
The one-tag branch satisfies the parallel context law:
$$
\mathsf{IntHistoryClassifier}(\Eone(\mathsf{append}(a,m)),\Eone(\mathsf{append}(a,n)))
\Longleftrightarrow
\mathsf{IntHistoryClassifier}(\Eone(m),\Eone(n)).
$$
The ten-field fit is the following. The traditional target is signed-integer
equality under a common unary magnitude context. The BEDC source consists of
$\mathsf{IntHistoryClassifier}$, $\mathsf{Unary}$, $\mathsf{append}$,
$\hsame$, and append cancellation. The carrier is the public signed-history
carrier. The classifier is the same-tag public classifier. Stability is common
unary prefix invariance. The ledger records the common visible unary prefix.
The core theorem is this iff statement. Exactness is preservation and
reflection of same-tag classifier equality. The standard bridge is ordinary
signed unary equality under equal magnitude shifts. The status is
$\mathsf{ProofTarget}$.
\end{theorem}
\begin{proof}
Apply Theorem~\ref{thm:int-history-classifier-same-tag-readback} to the
$\Ezero$ endpoints. With the unary closure of
$\mathsf{append}(a,m)$ and $\mathsf{append}(a,n)$, the left side reads back to
$\hsame(\mathsf{append}(a,m),\mathsf{append}(a,n))$ together with the two
unary payload witnesses. The right side reads back to
$\mathsf{Unary}(m)\land\mathsf{Unary}(n)\land\hsame(m,n)$.

For the zero-tag reverse direction, use congruence of $\mathsf{append}(a,-)$
on $\hsame(m,n)$ and unary append closure for the two payloads. Repacking by
the same-tag readback theorem gives the classifier on
$\Ezero(\mathsf{append}(a,m))$ and $\Ezero(\mathsf{append}(a,n))$.

For the zero-tag forward direction, read back the left classifier to obtain
$\hsame(\mathsf{append}(a,m),\mathsf{append}(a,n))$. Left append cancellation
under the common unary prefix $a$ recovers $\hsame(m,n)$, and the original
unary assumptions supply the two payload witnesses. Repacking by same-tag
readback gives $\mathsf{IntHistoryClassifier}(\Ezero(m),\Ezero(n))$.

The one-tag branch repeats the same two moves after replacing $\Ezero$ by
$\Eone$. The readback criterion is unchanged on payloads, so append congruence
proves the reverse implication, and left append cancellation proves the
forward implication before repacking the public classifier.
\end{proof}
\leanchecked{BEDC.Derived.IntUp.IntHistoryClassifier\_positive\_append\_context\_iff}

\begin{theorem}[Integer history classifier hsame transport]
\label{thm:int-history-classifier-hsame-transport}
For all $h,k,h',k':\mathsf{BHist}$, the public integer history classifier is
stable under replacement of both endpoints by $\hsame$-equal representatives:
$$
\hsame(h,h')\to
\hsame(k,k')\to
\mathsf{IntHistoryClassifier}(h,k)\to
\mathsf{IntHistoryClassifier}(h',k').
$$
The ten-field fit is the following. The traditional target is signed-integer
equality stable under representative replacement. The BEDC source consists of
$\mathsf{BHist}$, $\hsame$, $\Ezero$, $\Eone$, and
$\mathsf{IntHistoryClassifier}$. The carrier is
$\mathsf{IntHistoryCarrier}$. The classifier is the two-branch public
classifier. The stability endpoint is $\hsame$ transport. The ledger records
the visible sign tag together with the unary payload. The core theorem is this
transport statement. Exactness preservation is preservation of the classifier
under representative replacement. The standard bridge is replacement of
signed-unary representatives. The status is $\mathsf{ProofTarget}$.
\end{theorem}
\begin{proof}
Unfold the classifier witness for $\mathsf{IntHistoryClassifier}(h,k)$ and
split its visible signed branch.

In the zero-tag branch, keep the unary payload witnesses $m,n$, the payload
comparison $\hsame(m,n)$, and the stored mixed-tag exclusion. Compose the old
endpoint witness $\hsame(h,\Ezero(m))$ with the symmetry of $\hsame(h,h')$ to
obtain $\hsame(h',\Ezero(m))$. The same composition with $\hsame(k,k')$ gives
$\hsame(k',\Ezero(n))$.

It remains to rebuild the branch exclusion at the transported endpoints. Given
a hypothetical one-tag witness from $h'$ and $k'$, compose it back through
$\hsame(h,h')$ and $\hsame(k,k')$ to obtain the corresponding one-tag witness
from $h$ and $k$. The stored zero-branch exclusion then rules it out, so the
zero-tag classifier branch repacks at $h'$ and $k'$.

In the one-tag branch, keep the same unary payload data and payload
comparison. The supplied endpoint transports compose with the old witnesses to
produce $\hsame(h',\Eone(m))$ and $\hsame(k',\Eone(n))$.

For the transported exclusion, any zero-tag witness at $h'$ and $k'$ composes
back to a zero-tag witness at $h$ and $k$. The stored one-branch exclusion
rules it out, and the one-tag classifier branch repacks with the original
payload witnesses.
\end{proof}

\closureat{\IntUp}{\checkedCertStr}
